\section{The critical transition and its parity}\label{sec:critical}

The cubic threshold is visible in the normalized recurrence: the even ratio changes by order $m^2/N$ at step $m$, and summing this over $m$ steps produces $m^3/N$. At the transition, these increments cannot be discarded. They converge to a coupled differential system.

\begin{theorem}[Uniform critical approximation]\label{thm:critical}
For every fixed $T>0$ there is a finite $C_T$ such that, whenever $M\ge1$, $N\ge4M+2$ and $M^3/N\le T$,
\begin{align}
\max_{1\le m\le M}\left|R_{2m}-\cosh z_m\right|&\le C_T M^2/N,\label{eq:criticeven}\\
\max_{1\le m\le M}\left|R_{2m+1}-\shc z_m\right|&\le C_T M^2/N,\label{eq:criticodd}
\end{align}
where $z_m=(4\sqrt2/3)\sqrt{m^3/N}$. If $k_N^3/N\to c\in[0,\infty)$, then
\begin{align*}
R_{k_N}&\longrightarrow\cosh\left(\frac23\sqrt c\right)
&&\text{along even orders},\\
R_{k_N}&\longrightarrow\shc\left(\frac23\sqrt c\right)
&&\text{along odd orders at least three}.
\end{align*}
For $c>0$, the two limits are different.
\end{theorem}

\subsection{The exact two-component recursion}

Put $e_m=R_{2m}$ and $o_m=R_{2m+1}$ for $m\ge1$, with auxiliary seeds $e_0=o_0=1$. These seeds do not define $R_0$ or $R_1$. The coefficients multiplying $o_0$ vanish at the first step. Equation~\eqref{eq:momentrec} gives
\begin{equation}\label{eq:tworec}
e_m=A_me_{m-1}+B_mo_{m-1},\qquad
o_m=C_me_m+D_mo_{m-1},
\end{equation}
where
\begin{align*}
A_m&=\frac N{N-2m+1},&
B_m&=\frac{4(2m-1)(2m-2)}{3(N-2m+1)},\\
C_m&=\frac{3N}{(N-2m)(2m+1)},&
D_m&=\frac{(2m-2)N}{(N-2m)(2m+1)}.
\end{align*}
At $m=1$, these reproduce $R_2=N/(N-1)$ and $R_3=N^2/((N-1)(N-2))$. Under $N\ge4M+2$, all coefficients are nonnegative and
\begin{equation}\label{eq:rowsums}
A_m+B_m\le1+12m^2/N,\qquad
C_m+D_m\le1+4m^2/N,\qquad C_m\le2.
\end{equation}
For example, the first row sum has excess
\[
\frac{(2m-1)(8m-5)}{3(N-2m+1)}\le\frac{32m^2}{3N}.
\]
The second row sum is $N/(N-2m)$, with excess at most $4m/N$.

\subsection{A comparison system with an explicit solution}

For $x\ge0$ set
\[
z(x)=\frac{4\sqrt2}{3}\sqrt{x^3/N},\qquad
E(x)=\cosh z(x),\qquad O(x)=\shc z(x).
\]
Direct differentiation gives
\begin{equation}\label{eq:ode}
E'(x)=\frac{16x^2}{3N}O(x),\qquad
O'(x)=\frac3{2x}(E(x)-O(x)).
\end{equation}
The second identity extends continuously to $x=0$. The associated power series are
\[
E(x)=\sum_{\ell\ge0}\frac{(32x^3/(9N))^\ell}{(2\ell)!},\qquad
O(x)=\sum_{\ell\ge0}\frac{(32x^3/(9N))^\ell}{(2\ell+1)!}.
\]
For $x^3/N\le T$, termwise differentiation therefore gives
\begin{align}
|E(x)|+|O(x)|&\le K_T,&
|E(x)-O(x)|&\le K_Tx^3/N,\notag\\
|E'(x)|+|O'(x)|&\le K_Tx^2/N,&
|E''(x)|+|O''(x)|&\le K_Tx/N.\label{eq:derivatives}
\end{align}
The derivative estimates follow by factoring $x^2/N$ or $x/N$ from the differentiated series. The remaining series converge uniformly on the bounded interval for $x^3/N$.

\subsection{Local errors and stability}

Write $E_m=E(m)$ and $O_m=O(m)$. The residuals of \eqref{eq:tworec} are
\[
\rho_m^E=E_m-A_mE_{m-1}-B_mO_{m-1},\qquad
\rho_m^O=O_m-C_mE_m-D_mO_{m-1}.
\]
We claim that
\begin{equation}\label{eq:residual}
|\rho_m^E|+|\rho_m^O|\le K'_Tm/N.
\end{equation}
Taylor's theorem and \eqref{eq:derivatives} bound each backward difference minus its derivative at $m$ by $O_T(m/N)$. With $b_m=16m^2/(3N)$, the exact coefficients give
\[
A_m-1=O(m/N),\qquad B_m-b_m=O_T(m/N).
\]
For the second estimate, direct subtraction bounds the difference by a constant times $m/N+m^3/N^2$, which is $O_T(m/N)$. Now
\begin{align*}
\rho_m^E={}&E_m-E_{m-1}-E'(m)+b_m(O_m-O_{m-1})\\
&-(A_m-1)E_{m-1}-(B_m-b_m)O_{m-1}.
\end{align*}
Every term is $O_T(m/N)$. In particular, the second term is $O_T(m^4/N^2)$, and $m^3/N\le T$ makes it $O_T(m/N)$.

For the odd residual write
\[
c_m=\frac3{2m+1},\quad d_m=\frac{2m-2}{2m+1},\quad t_m=\frac N{N-2m}.
\]
Then $c_m+d_m=1$, $C_m=t_mc_m$, $D_m=t_md_m$ and $t_m-1=O(m/N)$. Removing $t_m$ changes the residual by $O_T(m/N)$. The remaining expression is
\[
d_m(O_m-O_{m-1})-c_m(E_m-O_m).
\]
After substituting \eqref{eq:ode}, the coefficient of $E_m-O_m$ is
\[
d_m\frac3{2m}-c_m=-\frac3{m(2m+1)}.
\]
Equation~\eqref{eq:derivatives} bounds its contribution by $O_T(m/N)$, proving \eqref{eq:residual}, including at $m=1$.

Let $H_m=\max(|e_m-E_m|,|o_m-O_m|)$ and $H_0=0$. Subtract the comparison equations from \eqref{eq:tworec}. First bound the even error, then use that bound in the odd equation. Positivity and \eqref{eq:rowsums} yield
\[
H_m\le(1+12m^2/N)(1+4m^2/N)H_{m-1}+3K'_Tm/N.
\]
The product of all multiplicative factors through step $M$ is at most
\[
\exp\left(\frac{16}N\sum_{m=1}^M m^2\right)\le e^{16T}.
\]
Summing the local errors proves $\max_{m\le M}H_m\le3K'_Te^{16T}M^2/N$, establishing \eqref{eq:criticeven} and \eqref{eq:criticodd}.

If $k^3/N$ stays bounded, then $m=O(N^{1/3})$ and the error tends to zero. For even $k=2m$, $z_m^2=4k^3/(9N)$. For odd $k=2m+1$, replacing $(2m)^3/N$ by $k^3/N$ changes it by $O(m^2/N)$. Continuity proves the stated limits.

\begin{corollary}[Small relative errors and sharpness]\label{cor:small}
If $k\to\infty$ and $k^3/N\to0$, then
\[
R_k-1\sim\frac{2k^3}{9N}\quad\text{for even }k,
\qquad
R_k-1\sim\frac{2k^3}{27N}\quad\text{for odd }k.
\]
Thus the constant $2/27$ in Theorem~\ref{thm:threshold} cannot be increased.
\end{corollary}
\begin{proof}
Take a fixed $T$, for instance $T=1$. The error $O(M^2/N)$ with $M$ equal to the relevant $m$ is smaller than $k^3/N$ by a factor of order $1/k$. Expand $\cosh z=1+z^2/2+O(z^4)$ and $\shc z=1+z^2/6+O(z^4)$.
\end{proof}

\begin{example}[A transition visible without a large formula]
When $k^3/N\to9/4$, the two limits are $\cosh(1)\approx1.5431$ and $\sinh(1)\approx1.1752$. Both exceed one. Their difference shows why an arbitrary sequence containing both parities need not have a common normalized limit.
\end{example}
